\documentclass[12pt]{article}
\usepackage[margin=0.72in]{geometry}
\usepackage{fontspec}
\setmainfont{Latin Modern Roman}
\usepackage{microtype}
\usepackage{fancyhdr}
\usepackage{titlesec}
\usepackage{enumitem}
\usepackage{needspace}
\usepackage{xcolor}
\usepackage[hidelinks]{hyperref}
\setlength{\parindent}{0pt}
\setlength{\parskip}{0.34em}
\setlength{\headheight}{15pt}
\titleformat{\section}{\large\bfseries\scshape}{}{0pt}{}
\titleformat{\subsection}{\normalsize\bfseries}{}{0pt}{}
\pagestyle{fancy}
\fancyhf{}
\fancyhead[L]{Whately Logic: Week 3 Test}
\fancyhead[R]{The Professor's Note}
\fancyfoot[C]{\thepage}
\renewcommand{\headrulewidth}{0.4pt}

\newcommand{\focusbox}[1]{%
\par\vspace{0.35em}\noindent\begingroup\setlength{\fboxsep}{6pt}\colorbox{gray!12}{\begin{minipage}{\dimexpr\linewidth-2\fboxsep\relax}#1\end{minipage}}\endgroup\par\vspace{0.45em}}
\newcommand{\studentline}{\par\vspace{0.60em}\hrule\vspace{0.64em}}
\newcommand{\shortline}{\par\vspace{0.38em}\hrule\vspace{0.48em}}

\newcommand{\twomeaningsblock}{%
\textbf{Two-faced word:}\shortline
\textbf{Meaning in the first use:}\shortline
\textbf{Meaning in the second use:}\shortline
\textbf{Fix the term to one clear meaning:}\studentline
\textbf{Standard form of the repaired claim (A/E/I/O):}\shortline
\textbf{Subject:}\shortline
\textbf{Predicate:}\shortline
\textbf{Quality:}\shortline
\textbf{Quantity:}\shortline
\textbf{After fixing the term and stating the precise form, does any conclusion the professor seems to draw follow? Explain.}\studentline}

\newcommand{\propblock}{%
\textbf{Standard form (A/E/I/O):}\shortline
\textbf{Subject:}\shortline
\textbf{Predicate:}\shortline
\textbf{Quality:}\shortline
\textbf{Quantity:}\shortline
\textbf{Does the precise form change what the claim appears to assert at first reading? Briefly explain.}\studentline}

\begin{document}
\begin{center}
{\LARGE\bfseries Week 3 Logic Test}\par\vspace{0.16em}
{\large\scshape The Case of the Professor's Note}\par\vspace{0.25em}
{A Whately Puzzle on Propositions, Quality, Quantity, and Shifting Terms}\par\vspace{0.45em}
\rule{0.82\textwidth}{0.4pt}
\end{center}

\section*{Name: \rule{2.35in}{0.4pt} \hfill Date: \rule{1.25in}{0.4pt}}

\focusbox{\textbf{Whately's guidance:} Turn ordinary language into a clear proposition before you judge it. Identify the subject and predicate. State the quality (affirmative or negative) and quantity (universal or particular). If a key word shifts meaning within the claim, fix the term to one sense first. Only then test whether a conclusion follows from the precise form. Logic examines the structure of what is asserted; it does not by itself settle biological facts.}

\section*{The Puzzle Story}

At the Academy of Straight Thinking the Biology Department occupies the west wing. After a session on animal classification, capabilities, and reproductive strategies, the professor leaves a long note on the main blackboard.

The note begins:

``Tell me where I'm mistaken.''

Then follows this list of ten statements drawn from the afternoon's discussion and the department's recent records:

\begin{enumerate}[leftmargin=*]
\item After reviewing wing structures and patterns of locomotion recorded in the department's extensive field notes from several habitats, it turns out that bats of every kind are equipped for flight.
\item In the examination of how birds move when they enter aquatic zones as illustrated by the lab's collection, the majority seem capable of soaring through water.
\item When one consults the full reproductive records maintained for living mammals in the academy archives, the conclusion is drawn that all such mammals give live birth to their young.
\item Cataloguing the external egg-laying practices of animals in the present set of specimens leads to the position that those which lay eggs belong only to birds or reptiles.
\item A thorough check of the anatomical features in the winged specimens on display reveals that there exist certain ones among them that lack any capacity for flight despite possessing those structures.
\item According to the classification of respiratory organs detailed in the reference volumes used by the department, animals equipped with gills fall entirely within the category of fish.
\item From the records of sensory capacities gathered for mammalian groups studied in the research program, most mammals have the capacity for echolocation.
\item Records kept of birds described as free in the academy's natural areas indicate that most free birds have free flight.
\item In the series of lectures on the aerial powers of mammals, it is asserted that true flight is true for bats.
\item A survey of flight abilities among the mammals included in the department's comparative anatomy studies shows the bat to be the only mammal that can achieve sustained flight.
\end{enumerate}

\clearpage

\section*{Part I: Translate Each Statement}

For each of the professor's ten statements, translate the claim into standard form. Identify the subject and predicate. State the quality (affirmative or negative) and quantity (universal or particular). 

For statements containing a word that may carry two different meanings, first identify the two-faced word, state both senses, fix the term to one clear meaning, and then provide the standard form.

\Needspace{0.30\textheight}
\subsection*{Statement 1}
\textbf{Claim:} ``After reviewing wing structures and patterns of locomotion recorded in the department's extensive field notes from several habitats, it turns out that bats of every kind are equipped for flight.''

\propblock

\Needspace{0.30\textheight}
\subsection*{Statement 2}
\textbf{Claim:} ``In the examination of how birds move when they enter aquatic zones as illustrated by the lab's collection, the majority seem capable of soaring through water.''

\propblock

\Needspace{0.30\textheight}
\subsection*{Statement 3}
\textbf{Claim:} ``When one consults the full reproductive records maintained for living mammals in the academy archives, the conclusion is drawn that all such mammals give live birth to their young.''

\twomeaningsblock

\Needspace{0.30\textheight}
\subsection*{Statement 4}
\textbf{Claim:} ``Cataloguing the external egg-laying practices of animals in the present set of specimens leads to the position that those which lay eggs belong only to birds or reptiles.''

\propblock

\Needspace{0.30\textheight}
\subsection*{Statement 5}
\textbf{Claim:} ``A thorough check of the anatomical features in the winged specimens on display reveals that there exist certain ones among them that lack any capacity for flight despite possessing those structures.''

\propblock

\Needspace{0.30\textheight}
\subsection*{Statement 6}
\textbf{Claim:} ``According to the classification of respiratory organs detailed in the reference volumes used by the department, animals equipped with gills fall entirely within the category of fish.''

\propblock

\Needspace{0.30\textheight}
\subsection*{Statement 7}
\textbf{Claim:} ``From the records of sensory capacities gathered for mammalian groups studied in the research program, most mammals have the capacity for echolocation.''

\propblock

\Needspace{0.30\textheight}
\subsection*{Statement 8}
\textbf{Claim:} ``Records kept of birds described as free in the academy's natural areas indicate that most free birds have free flight.''

\twomeaningsblock

\Needspace{0.30\textheight}
\subsection*{Statement 9}
\textbf{Claim:} ``In the series of lectures on the aerial powers of mammals, it is asserted that true flight is true for bats.''

\twomeaningsblock

\Needspace{0.30\textheight}
\subsection*{Statement 10}
\textbf{Claim:} ``A survey of flight abilities among the mammals included in the department's comparative anatomy studies shows the bat to be the only mammal that can achieve sustained flight.''

\propblock

\clearpage

\section*{Part II: Spiral from Week 1}

Take this claim drawn from one of the professor's statements:

\focusbox{``After a review of the department's records, most birds can soar through water, and therefore most birds in the collection should be moved to the aquatic tanks for further study.''}

\textbf{Conclusion:}\shortline

\textbf{Stated reason:}\shortline

\textbf{Hidden assumption:}\shortline

\textbf{Does the conclusion follow? Explain briefly, using the precise proposition form if it helps.}\studentline\studentline

\Needspace{0.38\textheight}
\section*{Part III: Spiral from Week 2}

Choose one statement from the professor's list (different from those used above) that does not obviously contain a shifting word. Rewrite it so that every important term has one clear meaning, then state the precise proposition form.

\textbf{Statement number chosen:}\rule{0.8in}{0.4pt}

\textbf{Rewritten version with all terms fixed:}\studentline\studentline

\textbf{Standard form after the rewrite:}\shortline

\textbf{Why was fixing the terms necessary before translating?}\studentline

\Needspace{0.42\textheight}
\section*{Part IV: What Can Logic Decide?}

Mark each item as \textbf{Logic can test this} or \textbf{Logic alone cannot decide this}. Then explain two of your choices in complete sentences.

\begin{enumerate}[leftmargin=*]
\item Whether a word such as ``free,'' ``live,'' or ``true'' is being used in two different senses in one of the professor's statements. \hfill \rule{1.6in}{0.4pt}
\item Whether most birds are in fact capable of soaring through water once the terms are fixed. \hfill \rule{1.6in}{0.4pt}
\item Whether every mammal gives live birth after ``live'' is fixed to a single meaning. \hfill \rule{1.6in}{0.4pt}
\item Whether a conclusion the professor draws from one of these claims actually follows from the precise proposition. \hfill \rule{1.6in}{0.4pt}
\item Whether bats are the only mammals capable of sustained flight as a matter of biological fact. \hfill \rule{1.6in}{0.4pt}
\item Whether the subject term in one of the claims is distributed. \hfill \rule{1.6in}{0.4pt}
\end{enumerate}

\textbf{Explain two choices:}\studentline\studentline\studentline

\Needspace{0.38\textheight}
\section*{Part V: Repair or Clarify One Claim}

Choose any one statement from the professor's list of ten. Rewrite it so that the proposition is stated in clear, standard language with no shifting terms. Then add one piece of real biological observation or evidence that would give the clarified claim stronger support.

\textbf{Statement number chosen:}\rule{0.8in}{0.4pt}

\textbf{Clear, repaired version:}\studentline\studentline

\textbf{One supporting observation or piece of evidence:}\studentline\studentline

\textbf{Why does the repaired version plus the added evidence make the claim easier to evaluate?}\studentline\studentline

\Needspace{0.28\textheight}
\section*{Final Whately Claim}

In one sentence, explain why a thinker must both fix any shifting terms and state the exact quantity and quality of a proposition before deciding whether the claim is useful or where the professor may be mistaken.

\studentline\studentline

\end{document}